% !TEX root =  Lecture13.tex


%%%%%%%%%%%%%%%%%%%%%%%%%%%%%%%%%%%%
% Recap/Agenda
%%%%%%%%%%%%%%%%%%%%%%%%%%%%%%%%%%%%
\begin{frame}
\frametitle{Module 2: Statistical Inference}
    \begin{itemize}
        \item \hl{Previously:} two \emph{categorical} variables: two-way tables, expected counts, and the $\chi^2$ statistic.
        \item \hl{This time:} back to \emph{numbers}, but now \hl{three or more groups} of them.
          \begin{itemize}
              \item why running every pairwise $t$-test is a trap;
              \item one statistic for $k$ groups: variability \emph{between} the groups over variability \emph{within} them, the $F$-statistic;
              \item its null distribution, the same two ways: shuffle the group labels, or use the $F$ model;
              \item the conditions, and which road survives when they fail;
              \item after a significant $F$: \emph{which} groups actually differ?
          \end{itemize}
        \item \hl{Reading:} IMS Chapter 22.
        \item \hl{Homework for today:}
        \item \hl{Homework for next class:}
    \end{itemize}
\end{frame}


%%%%%%%%%%%%%%%%%%%%%%%%%%%%%%%%%%%%
% Learning objectives:
%%%%%%%%%%%%%%%%%%%%%%%%%%%%%%%%%%%%
\begin{frame}
    \frametitle{Lecture Objectives}
    \goals{%
    \begin{itemize}
    \item explain why testing every pair inflates the false-positive rate, and what a single simultaneous test buys you;
    \item split the total variability into a \emph{between-group} and a \emph{within-group} part, and build $F$ out of them;
    \item obtain $F$'s null distribution by \emph{shuffling the group labels} \textbf{and} from the $F_{df_1, df_2}$ model, and explain why they agree;
    \item check the conditions (normality, and above all \emph{equal variances}), and know which road survives when they fail;
    \item after a significant $F$, say \emph{which} groups differ without walking back into the trap you started with.
    \end{itemize}}
    \objectives{%
    \begin{itemize}
    \item \textbf{(M2, LO3)} Analysis of Variance (ANOVA)
    \end{itemize}}
\end{frame}
